- Focus on one player, see how stats change of other players on the court with him
- 



player_stats.csv has Assists stat (AST)

Defining Clutch: 
The NBA's standard clutch situation is when the score is within 5 points during the final 5 minutes of the fourth quarter or overtime.

We can use the NBA's clutch endpoint to filter for those situations.
(LeagueDashPlayerClutch)

Factors: Pts made during Clutches

1. Clutch points — our primary measure
The logic is that players who consistently score during clutch situations may pose a greater offensive threat, causing opponents to devote more defensive attention to them.

Our potential proposition could be:

Proposition 4: Players with greater clutch scoring production tend to have higher offensive gravity.

One mathematical issue is that players have different amounts of clutch playing time.

\[
\text{Clutch PTS per 36 min}
=\frac{\text{Clutch PTS}}{\text{Clutch MIN}}\times36
\]

This would account for playing time, although we would need a minimum clutch-minutes requirement so extremely small samples don't dominate.

We could calculate:
\[
C_i=\frac{P_i}{M_i}\times36
\]

Where:
- \(C_i\) = clutch points per 36 minutes for player \(i\).
- \(P_i\) = total clutch points.
- \(M_i\) = total clutch minutes.

Adding a minimum threshold of clutch minutes of 30 to remove outliers.

Our mathematical definition becomes:
\[
C_i =
\begin{cases}
\dfrac{36P_i}{M_i}, & M_i\geq 30\\[6pt]
\text{undefined}, & M_i<30
\end{cases}
\]

Since some players might not have a clutch factor for our regression model. We will have to exlucde those players, hence limiting our sample size.





Players below the threshold will not receive a qualifying clutch scoring rate. We should never interpret their missing rate as zero.



Clutch assists is a future idea. Potential issue is statistical overlap --> multicollinearity.




-----------------------
Timeouts:

Seeing which players score and a timeout happens after/36 minutes

Proposition 5
Players with higher opposing timeout-response rates tend to have greater offensive gravity.

\[
\boxed{T_i=\frac{36N_i}{M_i}}
\]

Where:
- \(T_i\): Opposing timeout responses per 36 minutes for player \(i\).
- \(N_i\): Number of opposing timeouts occurring immediately after player \(i\)'s made field goals.
- \(M_i\): Total regular-season minutes played by player \(i\).